\changecolours{ScoutTeal}
\section{Prática \& Conclusão}

% ------------------------------------------------------------------------------
\begin{frame}[fragile]{Demonstração Prática: JSON vs 4 Tabelas Relacionais}
  \vspace{-0.2cm}
  \begin{columns}[t]
    \begin{column}{0.48\textwidth}
      \alert{Relacional: 4 Tabelas + 3 JOINs}\par
      \begin{lstlisting}[language=SQL]
SELECT c.nome, e.cidade, p.id_pedido, 
       i.produto, i.quantidade
FROM clientes c
JOIN enderecos e ON c.id = e.id_cliente
JOIN pedidos p ON c.id = p.id_cliente
JOIN itens_pedido i 
  ON p.id_pedido = i.id_pedido
WHERE c.id = 1;
      \end{lstlisting}
      \scriptsize Exige remontar objetos a partir de múltiplas linhas tabulares com redundâncias.
    \end{column}

    \begin{column}{0.48\textwidth}
      \alert{NoSQL: Documento Agregado}\par
      \begin{lstlisting}
{
  "_id": 1,
  "nome": "Maria Clara Silva",
  "endereco": {"cidade": "Taua", "uf": "CE"},
  "pedidos": [
    {
      "id_pedido": "PED-001",
      "itens": [
        {"produto": "Notebook", "qtd": 1}
      ]
    }
  ]
}
      \end{lstlisting}
      \scriptsize Leitura e escrita direta em 1 único I/O atômico de disco.
    \end{column}
  \end{columns}
\end{frame}

% ------------------------------------------------------------------------------
\begin{frame}{Matriz Comparativa: Relacional vs NoSQL}
  \vspace{-0.1cm}
  \resizebox{\textwidth}{!}{%
  \begin{tabular}{>{\raggedright\arraybackslash}p{3.5cm} >{\raggedright\arraybackslash}p{5.8cm} >{\raggedright\arraybackslash}p{6.0cm}}
    \toprule
    \textbf{Dimensão} & \textbf{Relacional (RDBMS)} & \textbf{Não-Relacional (NoSQL)} \\
    \midrule
    \textbf{Estrutura} & Esquema rígido e normalizado & Esquema dinâmico e flexível (\textit{schemaless}) \\
    \textbf{Escalabilidade} & Predominantemente Vertical (\textit{Scale-Up}) & Horizontal nativa em cluster (\textit{Scale-Out}) \\
    \textbf{Consistência} & Imediata e estrita (ACID) & Eventual ajustável (BASE / CAP) \\
    \textbf{Relacionamentos} & Otimizado para \texttt{JOIN}s locais & Desnormalização ou Grafos dedicados \\
    \textbf{Cenário Ideal} & Sistemas bancários, ERPs, sistemas financeiros & E-commerce, telemetria IoT, Big Data, IA / RAG \\
    \bottomrule
  \end{tabular}%
  }
\end{frame}

% ------------------------------------------------------------------------------
\twocolframe[title={Próximos Passos \& Atividades},leftcol=7cm,rightcol=7cm]{%
  \alert{Próximo Encontro: 06/10/2026}\par
  \itemseps[topsepi=2mm,itemsepi=3mm,labelsep=2mm,leftmargini=4mm]
  \textbf{Aula 02 -- Teorema CAP e BASE}
  \begin{itemize}
    \item O Teorema CAP de Eric Brewer detalhado.
    \item Teorema PACELC: a troca entre consistência e latência.
    \item Comparação formal: ACID vs BASE.
  \end{itemize}
}{%
  \alert{Tarefas do Estudante}\par
  \itemseps[topsepi=2mm,itemsepi=3mm,labelsep=2mm,leftmargini=4mm]
  \begin{itemize}
    \item Executar o script didático de apoio:
      \texttt{codigo/demonstracao\_aula01.py}.
    \item Iniciar os exercícios do Bloco 1 da \textbf{Lista de Exercícios 1 (AV3)}.
    \item Instalar o Docker para a oficina do primeiro sábado letivo (10/10).
  \end{itemize}
}
